\section{Related Work}
Web and computer-use agent benchmarks emphasize realistic execution. WebArena provides self-hosted web applications across e-commerce, forums, software development, and content management, with 812 end-to-end web tasks scored by functional correctness \href{https://arxiv.org/abs/2307.13854}{[WebArena]}. OSWorld moves this style to real desktop environments, with 369 tasks spanning web apps, desktop apps, file I/O, and multi-application workflows, each paired with setup and execution-based evaluation scripts \href{https://arxiv.org/abs/2404.07972}{[OSWorld]}. WorkArena and WorkArena++ focus on enterprise knowledge work in ServiceNow: WorkArena introduces 33 browser-based tasks, while WorkArena++ expands to 682 compositional tasks and ground-truth traces for planning and reasoning analysis \href{https://arxiv.org/abs/2403.07718}{[WorkArena]}, \href{https://arxiv.org/abs/2407.05291}{[WorkArena++]}. $\tau$-bench evaluates multi-turn tool-agent-user conversations in airline and retail domains using policy guidelines, API tools, final database-state comparison, and pass$^k$ reliability \href{https://arxiv.org/abs/2406.12045}{[$\tau$-bench]}. Across these benchmarks, the dominant validation pattern is executable or functional grading rather than mandatory double annotation of every task; DoneBench follows that tradition with reference replay, state-based grading, and task-quality audit, while treating human calibration as an optional strengthening layer.

These benchmarks share DoneBench's concern with state, tools, policies, and reproducible outcome grading, but they mostly keep the completion semantics inside the benchmark grader. DoneBench instead asks the agent to state those semantics before execution, scores atom-level agreement with hidden gold criteria, stress-tests the produced verifier on near-miss states, and then checks whether execution violates the agent's own criteria. This makes the tested object different from final task success: the central capability is knowing what would count as done before acting.

Evaluation frameworks and utility-assessment systems are adjacent but not substitutes. AgentEval automatically proposes task-specific utility criteria and applies them to LLM-powered applications \href{https://arxiv.org/abs/2402.09015}{[AgentEval]}, while DeepEval provides unit-test-like metrics, LLM-as-judge metrics, tracing, synthetic datasets, and CI workflows for LLM applications \href{https://deepeval.com/docs/introduction}{[DeepEval]}. These systems help practitioners define and run evaluations, but they do not package a stateful workflow benchmark where criterion construction itself is the model output under test.

Specification and verifier benchmarks provide another close boundary. CodeSpecBench evaluates executable behavioral specification generation for code, including function-level and repository-level Python specifications with correctness and completeness checks \href{https://arxiv.org/abs/2604.12268}{[CodeSpecBench]}. RESTestBench evaluates LLM-generated REST API tests from natural-language requirements using three REST services, precise and vague requirement variants, and requirement-based mutation testing \href{https://arxiv.org/abs/2604.25862}{[RESTestBench]}. VerifyLLM verifies robotic task plans before execution by translating natural-language instructions into Linear Temporal Logic and analyzing action sequences \href{https://arxiv.org/abs/2507.05118}{[VerifyLLM]}. DoneBench borrows the insight that executable specifications and adversarial invalid cases are essential, but shifts the domain from program or robot-plan verification to everyday tool-using workflow agents and ties specification grounding to downstream execution and self-violation metrics.
